\scanpage{27}
% archive sheet number: 26
Soit $K$ un corps top. isomorphe à $\mathbb C$ [isom non choisi !]

Alors on a le foncteur \inserted{mult.} $\unsure{\xi|\zeta}=\xi_K$ sur la catégorie des motifs sur $K$, à valeurs en des $\mathbb Z$-modules, et des isom
\[
\begin{aligned}
 \xi\otimes_{\mathbb Z}\mathbb Z_{\unsure{\ell|l}} &\overset{\alpha_\ell}{\simeq} \xi_\ell\\
 \xi\otimes_{\mathbb Z}K &\overset{\alpha_\infty}{\simeq} \xi_K
\end{aligned}
\]
Considérons
\[
 \xi(H^2(\mathbb P^1,\deleted{\unsure{\mathbb Q}}\mathbb Z(0)))=\xi(\mathbb Z(-1))
\]
c'est isomorphe à $\mathbb Z$, l'isom n'étant pas choisi. Le choix d'un tel isom \inserted{\unsure{i.e.|ni} d'un des gén.''} revient à celui d'une \uline{orientation} sur $\mathbb P^1_K$ \unsure{[|(}qui est de façon naturelle une variété \deleted{analytique} réelle, isomorphe à $S^2$), ou ce qui revient au \unsure{$=$|même}, au choix de l'une des deux structures \unsure{complexes|orientées} sur $\mathbb P^1_K$, provenant du \deleted{\unsure{2}} \inserted{choix d'un} isom $K\simeq\mathbb C$, i.e. le choix d'\uline{\unsure{une|un}} des racines de $-1$ dans $K$. Comme $\pi\in\mathbb R\subset\unsure{K|\mathbb R}$ est bien déterminé comme

